\vfill\pagebreak
\section{Repeating while a condition holds: \texttt{while}}
\label{sec:control:while}

Many loops have the same shape: test a condition at the top, break if it does
not hold, otherwise do the work. \token{while} writes that shape directly.
\token{while condition \{ ... \}} tests \token{condition} before each iteration,
runs the block if it is \token{true}, and leaves the loop as soon as it is
\token{false}.

The next listing counts the digits of a number by dividing it by ten until a
single digit is left. Figure~\ref{fig:control:while} draws its flow: the arrow
from the bottom of the block goes back up to the test, not to the start of the
block.

\begin{lstlisting}[style=coloredverbatim, caption=Counting digits with a \token{while} loop, label=lst:control:while]
use std::io;

fn main() {
  let mut n = read!i32(ask-> "A positive number: ");
  let mut digits = 1;
  while n >= 10 {
    n = n / 10;
    digits += 1;
  }
  println("That number has ", digits, " digits.");
}
\end{lstlisting}
\vspace{-10pt}%
\begin{lstlisting}[style=bashVerb, escapechar=@+]
@+\textcolor{teal!80}{alice@dev:\mtilde}@+$ ./main
A positive number: 40953
That number has 5 digits.
\end{lstlisting}

\begin{figure}[h]
  \centering
  \begin{adjustbox}{max width=\linewidth}
    \begin{tikzpicture}[node distance=6mm]

      \node[cfterm] (start) {start of main};
      \node[cfstep, below=of start] (init) {let mut n = read!i32(...);\\let mut digits = 1;};
      \node[cftest, below=of init] (test) {n >= 10};
      \node[cfstep, below=8mm of test] (body) {n = n / 10;\\digits += 1;};
      \node[cfstep, below=of body] (print) {println(..., digits, ...);};
      \node[cfterm, below=of print] (end) {end of main};

      \draw[cfflow] (start) -- (init);
      \draw[cfflow] (init) -- (test);
      \draw[cfflow] (test) -- node[cflab, right] {true} (body);
      \draw[cfflow] (body.east) -- ++(14mm, 0) |- (test.east);
      \draw[cfflow] (test.west) -- node[cflab, above] {false} ++(-14mm, 0) |- (print.west);
      \draw[cfflow] (print) -- (end);
    \end{tikzpicture}
  \end{adjustbox}
  \caption{\label{fig:control:while} The \token{while} loop of
    Listing~\ref{lst:control:while}. The condition is tested before every
    iteration, including the first.}
\end{figure}


Because the test comes first, a \token{while} can run its block zero times: with
\token{7} as input, \token{n >= 10} is \token{false} from the start, and the
program reports a single digit without entering the loop. The same program written with
\token{loop} needs the test and the \token{break} spelled out:

\begin{lstlisting}[style=coloredverbatim, caption=The same loop without \token{while}]
use std::io;

fn main() {
  let mut n = read!i32(ask-> "A positive number: ");
  let mut digits = 1;
  loop {
    if !(n >= 10) {
      break;
    }
    n = n / 10;
    digits += 1;
  }
  println("That number has ", digits, " digits.");
}
\end{lstlisting}

\token{break} and \token{continue} work in a \token{while} exactly as in a
\token{loop}. A \token{while}, however, has no value: it may stop because its
condition became false, and then no \token{break} has provided one. So
\token{break} inside a \token{while} never carries a value.

\subsection{Choosing between \texttt{while} and \texttt{loop}}

A \token{while} whose condition is always \token{true} is a \token{loop} under
another name, and the compiler insists on the real name. \token{while true \{}
is rejected with ``infinite loop with always true test is not allowed, rather use
a loop construct''.

\begin{lstlisting}[style=coloredverbatimError, escapechar=@, caption=\token{while true} is written \token{loop}]
use std::io;

fn main() {
  let mut total = 0;
  while @\hb{true}@ {
    let n = read!i32(ask-> "A number (0 to stop): ");
    if n == 0 {
      break;
    }
    total += n;
  }
  println("Total: ", total);
}
\end{lstlisting}

Which of the two fits depends on where the loop decides to stop. When the
decision can be taken before the work, as in Listing~\ref{lst:control:while},
use \token{while}. When the work has to happen first -- the number must be read
before it can be compared with zero, the month must be asked before it can be
checked -- the decision sits in the middle or at the end of the block, and the
loop is a \token{loop} with a \token{break} there.

\notebox{Ymir has no \textit{do-while} loop, which some languages provide to test
  the condition after the block rather than before it. The \token{loop} covers
  that case, with the test written at the end of the block:
  \token{loop \{ ... if !condition \{ break; \} \}}.}
